%In the category of varieties the morphisms between objects have a special name

To appease the critical mind of the reader, a few more concrete examples are presented to give a feeling for the previous abstract definitions. 


Maps between geometric spaces can tell us how they relate to each other, or simplify 
the derivation of properties for one space if they can be inherited by the existence of a map to another space,
already satisfying them. In case of varieties the maps between them satisfy a great deal of structure.
\begin{definition}[Morphisms between varieties]
    Let \( X, Y \) be affine varieties, then a continuous map \( f : X \to Y \) is regular
    if it is morphism of the corresponding affine schemes.
\end{definition}

\begin{example}
    Let \( X = \{ (x_1, x_2) \in \mathbb{R}^2 | x_1^2 - x_2 = 0 \} \), \( Y = \{ (y_1, y_2) \in \mathbb{R}^2 | y_1 - 3 y_2 = 0 \} \)
    and \( f : X \to Y, (x_1, x_2) \mapsto (3 x_1^2, x_2) \).
    Clearly, the defining polynomials are irreducible and consequentially their solution sets as well, making them affine varieties.
    To verify that there is a morphism on the scheme construction, it is helpful too look at the map that \( f \) induces on the coordinate rings,
    which are 
    \[ 
        k[X] = \mathbb{R}[x_1, x_2] / \langle x_1^2 - x_2 \rangle \simeq \mathbb{R}[x_1] \oplus \mathbb{R} 
    \]
    and 
    \[ 
        k[Y] = \mathbb{R}[y_1, y_2] / \langle y_1 - 3 y_2 \rangle \simeq \mathbb{R}[y_1].
    \]
    Consider the map \( f^\ast : k[Y] \to k[X],\ g(y_1, y_2) \mapsto g(3x_1^2, x_2) \equiv g \mapsto g \circ f \), for which it is readily 
    checked that this is even a homomorphism of rings. The only subtlety is that one works over quotient rings, therefore equivalent polynomials
    should be sent to the same equivalence class in \( k[X] \), therefore suppose 
    \( h, g \in \mathbb{R}[y_1, y_2] \setminus \langle y_1 - 3y_2 \rangle,\ r \in \mathbb{R}[y_1, y_2] \) are s.t. \( h - g = r(y_1 - 3 y_2) \),
    then \( f^\ast(h - g) = f^\ast(r)(f^\ast(y_1) - 3 f^\ast(y_2)) = f^\ast(r)(3 x_1^2 - 3 x_2) = 3f^\ast(r)(x_1^2 - x_2) = 0 \).
    \newline
    The spectra of each variety are given by \( \spec{k[X]} \simeq \spec{\mathbb{R}[x_1]} \sqcup \spec{\mathbb{R}} \simeq \spec{\mathbb{R}[x_1]} \)
    and \( \spec{k[Y]} = \spec{\mathbb{R}[y_1]} \).
    \( f \) induces on them the map 
    \( f^\# : \spec{\mathbb{R}[x_1]} \mapsto \spec{\mathbb{R}[y_1]} , \mathfrak{m} \mapsto \left(f^{\ast}\right)^{-1}(\mathfrak{m}) \)
    (since the preimage of any prime ideal is again a prime ideal under a homomorphism).
    Now take \( g \in k[Y] \) and the associated affine open \( D_g = \{ \mathfrak{p} \in \spec{k[Y]} | g \notin \mathfrak{p} \} \),
    then for \( h \in \mathcal{O}(D_g) = \{ \frac{r}{g^n} | r \in k[Y] \} \), we extend \( f^\ast \) on \( \mathcal{O}(D_g) \)
    via \( f^\ast(h) = f^\ast(r) \frac{1}{f^\ast(g)^n} \), which is well defined in \( \mathcal{O}(D_{f^\ast(g)}) \) in \( \mathcal{O}(k[X]) \).
    By construction each \( \mathfrak{p} \in \mathcal{O}(D_{f^\ast(g)}) \) is sent by \( f^\# \) precisely to some \( \mathcal{O}(D_g) \).
    That this plays well with the restrictions is again a consequence of \( f^\ast \) being a homomorphism.
    \newline
    \newline
    Finally we see that the actual morphism of sheaves is \( f^\# : \spec{k[X]} \to \spec{k[Y]} \), but to really understand it, one needs the associated map
    \( f^\ast : k[Y] \to k[X] \) on the coordinate rings. The confusing part to this, is the contravariance of the \( \text{Spec } \)-functor. Put succinctly 
    one can show that \( \text{Hom}_{\text{Sh}}(\spec{B}, \spec{A}) \simeq \text{Hom}_{\text{Ring}}(A, B) \), so morphisms on sheaves/varieties induce 
    maps in the opposite directions of the associated sections/coordinate rings. As in our case \( A, B \) are typically quotient rings of \( \mathbb{K}[x_1, \dots, x_n] \),
    which already implies that the maps are polynomial in each coordinate.
\end{example}

\begin{example}[Alg. Geo. is not limited by polynomials]
    Consider the generalized Pareto distribution \( \text{GPD}(\mu, \sigma, \xi) \), with CDF/PDF
    \begin{align*}
        F_{(\mu,\sigma,\xi)}(x) &= \begin{cases}
1 - \left(1+ \frac{\xi(x-\mu)}{\sigma}\right)^{-1/\xi} & \text{for }\xi \neq 0, \\
1 - \exp \left(-\frac{x-\mu}{\sigma}\right) & \text{for }\xi = 0,
\end{cases} \\
            f_{(\mu,\sigma,\xi)}(x) &= \frac{1}{\sigma}\left(1 + \frac{\xi (x-\mu)}{\sigma}\right)^{- \left(1 + 1/\xi\right)}
    \end{align*}
    \[
    \]
    with \( \mu, \xi \in \mathbb{R}, \sigma \in (0, \infty) \) and \( x \geq \mu \) if \( \xi \geq 0 \) otherwise \( \mu \leq x \leq \mu - \frac{\sigma}{\xi} \). 
    Its first 2 moments are given for \( \xi < \frac{1}{2} \) by
    \begin{align} \label{momentsofgpd}
        \expv{X} &= \mu + \frac{\sigma}{1 - \xi} \\
        \var{X} &= \frac{\sigma^2}{(1 - \xi)^2 (1 - 2\xi)}.
    \end{align}
    \newline
    Suppose we are given \( \expv{X}, \var{X} \) and should find all admissable \( \mu, \xi, \sigma \) (akin to the method of moments). Ignoring for now
    that \( \sigma \geq 0 \) (which can be incorporated after introducing semi-algebraic sets \ref{semialgset}) Equation \ref{momentsofgpd}
    are a set of equations with rational polynomials. Transforming this into
    \begin{align*} 
        0  &= \mu\left(1 - \xi \right) + \sigma - \expv{X}\left(1 - \xi \right) \\
        0  &= \sigma^2 - \var{X}(1 - \xi)^2 (1 - 2\xi) \\
        0  &= (1 - \xi) (1 - 2\xi)y - 1
    \end{align*}
    defines a variety in \( \aff{\mathbb{R}}{4} \). The reason for the last equation is that we do not admit
    solutions, where rational polynomials blow up, hence at all admissable points there exist an inverse to the denominator, 
    represented by the dummy variable \( y \). This trick is applicable to all rational polynomials, although the last equation
    is in general the product of all denominators (which would have been \(0 = (1 - \xi)^3 (1 - 2\xi)y - 1 \) but since reducible polynomials
    exactly vanish if their irreducible factors do, we can neglect the exponent).
    \newline
    \newline
    Now supose we are given \( \xi = \frac{6}{5} \) (close to the factor for the famous Pareto principle)
    and \( 2 \) points of the empirical distribution function \( \{(x_0, p_0), (x_1, p_1)\} \) and should find suitable \( \mu, \sigma \) 
    (this is for purely demonstrative purposes, the author recommends to abstain from such statistics in general).
    The corresponding equations are:
    \begin{align} \label{gpdfitprob}
        p_0 &= 1 - \left(1+ \frac{\frac{6}{5}(x_0-\mu)}{\sigma}\right)^{-5/6} \\ 
        p_1 &= 1 - \left(1+ \frac{\frac{6}{5}(x_1-\mu)}{\sigma}\right)^{-5/6}.
    \end{align}
        Now define two new variables \( z_1, z_2 \), which satisfy the following equations 
    \begin{align} \label{rootutilityvars}
        z_1^{6} &= 1 + \frac{ \frac{6}{5} (x_0-\mu)}{\sigma} \\
        z_2^{6} &= 1 + \frac{ \frac{6}{5} (x_1-\mu)}{\sigma}
    \end{align}
    Using the trick from above, this is equivalent to the following polynomial system
    \begin{align*}
        \sigma z_1^{6} &= \sigma + \frac{6}{5} (x_0-\mu)\\ 
        \sigma z_2^{6} &= \sigma + \frac{6}{5} (x_1-\mu) \\
        \sigma y - 1 &= 0.
    \end{align*}
    Now reexpress the original system from (\ref{gpdfitprob}) via 
    \begin{align*}
        z_1^{5}(p_0 - 1) + 1 &= 0 \\ 
        z_2^{5}(p_0 - 1) + 1 &= 0.
    \end{align*}
    We have successfully translated our original problem into an algebraic variety.
    Additionally, the last equations for \( z_1, x_2 \) are straightforward to solve, which in turn
    resolves the equations for \( \mu, \sigma \) in \ref{rootutilityvars}.
    To review what was done, a general equation \( f(x_1, \dots, x_n)^\frac{a}{b} + g(x_1, \dots, x_n) = 0 \), where \( f, g \) are polynomials
    with exponent in \( \mathbb{Q}_{>0} \) can be turned into a system of polynomials without fractional exponents,
    by defining new variables \( z^b = f(x_1, \dots, x_n) \) and rewriting the original equation to \( z^a + g(y_1, \dots, y_n) = 0 \), this can be done for each rational exponent,
    but often one can save some coordinates, by reducing first the number of terms with such exponents.
    \newline
    \newline
    Taking a step back, the thoughtful reader might notice, that this defines a morphism between schemes. Let 
    \[ X = \{ (x_1, \dots, x_n, z) \in \mathbb{R}^{n+1} | z_1^b - f(x_1, \dots, x_n) = 0 \} \] 
    \[ Y = \{ (y_1, \dots, y_n, z) \in \mathbb{R}^{n+1} | z_2^a + g(y_1, \dots, y_n) = 0 \} \]
    the corresponding morphism of coordinate rings \( k[Y] \to k[X] \) is not exactly trivial, but will send \( (y_1, \dots, y_n) \) to \( (x_1, \dots, x_n) \) 
    and \( z_2 \) to some polynomial \( h(z_1, x_1, \dots, x_n) \) s.t. \( h(z_1, x_1, \dots, x_n) + f(x_1, \dots, x_n) = 0 \).
    Further, there is even a geometric interpretation. 
    Over \( \mathbb{C} \), \( z_1^b - f(x_1, \dots, x_n) = 0  \) for fixed \( (x_1, \dots, x_n) \in \mathbb{C}^n \setminus V(\langle f \rangle) \), has precisely \( b \)
    solutions for \( z_1 \), similarly \( z_2^a + g(x_1, \dots, x_n) \) has \( a \) such roots, giving \( a \cdot b \) solutions in total.
    However, for \( (x_1, \dots, x_n) \in V(\langle f \rangle + \langle g \rangle) \), there is only one solution, namely \( z_i = 0 \) for \( i \in \{1, 2\} \).
    The schemes \( X, Y \) build a cover of the curves \( f(x_1, \dots, x_n), g(x_1, \dots, x_n) \) with \( a \), \( b \) respectively, branches for each point,
    except at points in the vanishing sets of \( f, g \), where they compress down to a point. This phenomenon is referred to as tame ramification.
\end{example}
